\chapter{Prepare Your Dev Environment}
\label{chap:fpgasteel-devenv}

\section{Introduction}

The development host is the \term{PC} on which \prjname{FPGAsteel}'s \gterm{fpga}{FPGA} project is compiled and the bare-metal applications that run on the \gterm{hps}{HPS} are cross-compiled; it is also the machine the \term{DE1-SoC} board stays connected to while it runs. This chapter gives it the setup every later chapter assumes: it installs the tools this part is built and debugged with, makes them reachable from the command line, and grants access to the two \term{USB} devices the board is reached through, the \gterm{usbblaster}{USB-Blaster II}\footnote{\glsentrydesc{usbblaster}} and the serial console.

Unlike \prjname{FPGAlix}, it is not possible here to debug the software over the network, since by design this project has no operating system. The only way to observe it from the host is therefore the serial console's output, together with the \gterm{usbblaster}{USB-Blaster II}'s \gterm{jtag}{JTAG}\footnote{\glsentrydesc{jtag}} interface, used to inspect and debug both the hardware and the software.\footnote{Developing inside a \gterm{vm}{Virtual Machine} is discouraged for this: forwarding the \gterm{usbblaster}{USB-Blaster II} through a hypervisor's \term{USB} redirection is a long-standing, widely reported source of \gterm{jtag}{JTAG} failures. Developing directly on physical \term{Ubuntu 22.04} hardware is what is documented and known to work.}

\section{Host requirements}

\prjname{FPGAsteel} requires an \term{x86\_64} machine running \term{Ubuntu 22.04 LTS}, with the four pieces of software covered in the following sections installed and working: \gterm{quartusprime}{Quartus Prime}, the \gterm{armbaremetaltoolchain}{\term{ARM} bare-metal toolchain}, \gterm{openocd}{OpenOCD}\footnote{\glsentrydesc{openocd}}, and \gterm{eclipseembeddedcdt}{Eclipse Embedded CDT}\footnote{\glsentrydesc{eclipseembeddedcdt}}.

\section{Required software}

\subsection{apt packages}
\label{subsec:fpgasteel-apt-packages}

The first step is installing a handful of additional \term{Ubuntu} system programs, commonly called apt packages, through \term{Ubuntu}'s own package manager. These aren't specific to this project, but the rest of this chapter, and the project as a whole, relies on them being in place.

\begin{shellcodehost}
sudo apt install -y build-essential bison flex libssl-dev libncurses-dev u-boot-tools git wget curl openjdk-17-jre python3 python3-numpy python3-matplotlib dosfstools picocom pkg-config libopencv-dev
\end{shellcodehost}

\begin{itemize} 
    \item |build-essential|, |bison|, |flex|, |libssl-dev|, |libncurses-dev|, |u-boot-tools|, |git|: building \repopath{repos/u-boot-socfpga} and its boot scripts.
    \item |wget|, |curl|: fetching files over \term{HTTP}/\term{FTP} and their secure variants, from the command line.
    \item |dosfstools|: provides \cmdpath{mkfs.vfat}, used to format the \term{SD card}.
    \item |picocom|: terminal emulator for the board's serial console.
    \item |openjdk-17-jre|: needed by \term{Eclipse}, installed later in this chapter (Paragraph~\ref{subsec:fpgasteel-install-paths}).
    \item |python3|, |python3-numpy|, |python3-matplotlib|: used by \cmd{cv_bsp_generator.py} (Section~\ref{subsec:fpgasteel-bsp-headers}) and \cmd{debugTools/view_frame.py} (Section~\ref{sec:fpgasteel-cache-forwarding-debug-tooling}).
    \item |pkg-config|, |libopencv-dev|: needed to build \cmd{debugTools/linux_dualcore_emulator}, the host-side emulator described in Section~\ref{sec:fpgasteel-dualcore-dip-debug-tooling}.
\end{itemize}

The \term{ARM} toolchain and \gterm{openocd}{OpenOCD} need no |apt| packages: both are prebuilt \gterm{xpackbuilds}{\term{xPack}} tarballs bundling their own shared libraries. \term{Quartus} needs none either: its binaries are \term{64-bit} and it ships its own private \term{JRE}.

\subsection{Install the software, and where it lives}
\label{subsec:fpgasteel-install-paths}

\begin{table}[H]
\centering
\begin{tabularx}{\textwidth}{@{} p{3.4cm} l X @{}}
\toprule
\textbf{Software} & \textbf{Version} & \textbf{Install path} \\
\midrule
\gterm{quartusprime}{Quartus Prime} & \texttt{22.1} & \texttt{\seqsplit{\textasciitilde{}/Programs/Altera/22.1/}} \\
\lightmidrule
\term{ARM} toolchain (\term{Embedded GCC}) & \texttt{15.2.1-1.1} & \texttt{\seqsplit{\textasciitilde{}/Programs/Arm/toolchain/xpack-arm-none-eabi-gcc-15.2.1-1.1/}} \\
\lightmidrule
\gterm{openocd}{OpenOCD} & \texttt{0.12.0-7} & \texttt{\seqsplit{\textasciitilde{}/Programs/openocd/xpack-openocd-0.12.0-7/}} \\
\lightmidrule
\term{Eclipse IDE for Embedded C/C++} & \texttt{2026-03 R} & \texttt{\seqsplit{\textasciitilde{}/Programs/Eclipse/embedcpp-2026-03-R-linux-gtk-x86\_64/}} \\
\bottomrule
\end{tabularx}
\caption{Required software and installation paths.}
\label{tab:fpgasteel-dev-software}
\end{table}

Each of these four programs should be installed at the path shown in Table~\ref{tab:fpgasteel-dev-software} above; the user's home folder is sufficient for all of them, a system-wide installation is not required. Copies of all four installers are also available from \nameref{chap:fpgasteel-downloads}.

\subsection{Environment variables}
\label{subsec:fpgasteel-env-vars}

Add to {\catcode`\~=12 \extpath{~/.bashrc}}, matching the install paths from Paragraph~\ref{subsec:fpgasteel-install-paths}:

\begin{shellcodehost}
export QUARTUS_ROOTDIR="$HOME/Programs/Altera/22.1/quartus/"
export QSYS_ROOTDIR="${QUARTUS_ROOTDIR}/sopc_builder/bin"
export PATH="${QSYS_ROOTDIR}:$PATH"
export PATH="${QUARTUS_ROOTDIR}/bin:$PATH"
export PATH="$HOME/Programs/Arm/toolchain/xpack-arm-none-eabi-gcc-15.2.1-1.1/bin:$PATH"
export PATH="$HOME/Programs/openocd/xpack-openocd-0.12.0-7/bin:$PATH"
export PATH="$HOME/Programs/Eclipse/embedcpp-2026-03-R-linux-gtk-x86_64:$PATH"
\end{shellcodehost}

Each line prepends to \term{PATH} rather than appending, and in this specific order, on purpose: \cmdpath{openocd} and \cmdpath{arm-none-eabi-gcc} both exist as regular \term{Ubuntu} packages, and \term{Quartus 22.1} additionally bundles its own, older \gterm{openocd}{OpenOCD} (0.11.0) at \extpath{quartus/bin/openocd}. Appending would only fix the ordering among these lines themselves; anything already on \term{PATH} before this block runs, including \extpath{/usr/bin}, would still be found first. Prepending, instead, guarantees the \term{xPack} builds this chapter installs always win, over both \term{Quartus}'s bundled \gterm{openocd}{OpenOCD} and any distro package of the same name, whether or not either is installed. Since the last line to run ends up first on \term{PATH}, the order above puts the \term{xPack OpenOCD} ahead of \lstinline[style=inlineFsExt]!${QUARTUS_ROOTDIR}/bin!.

\term{QUARTUS\_ROOTDIR} (or, as a fallback, \term{QSYS\_ROOTDIR}) is not just a convenience variable: \gterm{openocd}{OpenOCD}'s own \gterm{jtag}{JTAG} config script (\repopath{altera-usb-blaster2.cfg}, used from \term{Eclipse}'s \term{Debug Configuration}, Paragraph~\ref{subsec:fpgasteel-hello-debug-config}) reads it to locate \extpath{blaster_6810.hex}, the \gterm{usbblaster}{USB-Blaster II} firmware file it needs to talk to the \gterm{jtag}{JTAG} adapter.

Reload the updated \gterm{environmentvariables}{environment variables}\footnote{\glsentrydesc{environmentvariables}} with the following command, or open a new shell:

\begin{shellcodehost}
source ~/.bashrc
\end{shellcodehost}

\subsection{USB-Blaster II and serial console access}
\label{subsec:fpgasteel-usb-blaster}

Beyond the packages installed in Paragraph~\ref{subsec:fpgasteel-apt-packages}, the board itself requires access to two \term{USB} devices: the \gterm{usbblaster}{USB-Blaster II}, used for \gterm{jtag}{JTAG} programming and debugging, and the \term{UART-to-USB} serial console. By default, \gterm{udev}{udev}\footnote{\glsentrydesc{udev}} restricts both device nodes to root, which would force every \gterm{openocd}{OpenOCD} invocation and every serial connection to run as root. The following rule grants the invoking user's group access to the \gterm{usbblaster}{USB-Blaster II} instead:

\begin{shellcodehost}
sudo tee /etc/udev/rules.d/51-usbblaster.rules << 'EOF'
SUBSYSTEM=="usb", ATTR{idVendor}=="09fb", ATTR{idProduct}=="6001", MODE="0666", GROUP="plugdev"
SUBSYSTEM=="usb", ATTR{idVendor}=="09fb", ATTR{idProduct}=="6002", MODE="0666", GROUP="plugdev"
SUBSYSTEM=="usb", ATTR{idVendor}=="09fb", ATTR{idProduct}=="6003", MODE="0666", GROUP="plugdev"
SUBSYSTEM=="usb", ATTR{idVendor}=="09fb", ATTR{idProduct}=="6010", MODE="0666", GROUP="plugdev"
SUBSYSTEM=="usb", ATTR{idVendor}=="09fb", ATTR{idProduct}=="6810", MODE="0666", GROUP="plugdev"
EOF
sudo udevadm control --reload-rules
sudo udevadm trigger
sudo usermod -aG plugdev "$USER"
\end{shellcodehost}

The serial console needs no equivalent custom rule. The board's \term{UART-to-USB} connector enumerates as \extpath{/dev/ttyUSB0}, used at 115200 baud, 8-N-1, and \term{Ubuntu}'s own default rule already assigns every serial device to the \sysid{dialout} group. Add the user to that group too:

\begin{shellcodehost}
sudo usermod -aG dialout "$USER"
\end{shellcodehost}

Log out and back in for both group changes to take effect.\footnote{Or run \cmdpath{newgrp plugdev} / \cmdpath{newgrp dialout} on the already-open console to apply the change immediately, without logging out.}

\subsection{Editor and debugger roles}
\label{subsec:fpgasteel-editor-roles}

\term{Eclipse} is used exclusively to debug over \gterm{jtag}{JTAG}/\gterm{openocd}{OpenOCD}: it never compiles anything, and it is not used to write \gterm{vhdl}{VHDL} either. Both the \gterm{vhdl}{VHDL} sources under \repopath{hw/} and the \term{C/C++} sources under \repopath{sw/} are written in \gterm{visualstudiocode}{Visual Studio Code}\footnote{\glsentrydesc{visualstudiocode}}, with the following extensions:

{\footnotesize
\begin{itemize}
    \item \gterm{vhdl}{VHDL}/\term{Verilog} linting, snippets, project awareness --- \urlpath{teros-technology.teroshdl} (\term{TerosHDL}).
    \item \gterm{vhdl}{VHDL} syntax highlighting and language support --- \urlpath{rjyoung.vscode-modern-vhdl-support}.
    \item \term{C/C++} \term{IntelliSense}, navigation --- \urlpath{ms-vscode.cpptools} (+ \urlpath{ms-vscode.cpptools-extension-pack}).
    \item Makefile awareness (every \repopath{sw/} project builds via \gterm{make}{\cmdpath{make}}, Section~\ref{sec:fpgasteel-hello-building}) --- \urlpath{ms-vscode.makefile-tools}.
    \item Editing the \term{Python} debug tools shipped with some examples (e.g. \cmd{debugTools/view_frame.py}, Section~\ref{sec:fpgasteel-cache-forwarding-debug-tooling}) --- \urlpath{ms-python.python} (+ \term{Pylance}, \term{debugpy}).
    \item Viewing the datasheet/manual \term{PDF}s referenced throughout \repopath{docs/} without leaving the editor --- \urlpath{tomoki1207.pdf}.
\end{itemize}
}

